\documentclass[11pt]{article}
\usepackage{amsmath,amssymb,amsthm}
\usepackage[margin=1in]{geometry}
\usepackage{hyperref}
\usepackage{xurl}
\emergencystretch=3em

\newtheorem{theorem}{Theorem}
\newtheorem{lemma}[theorem]{Lemma}
\newtheorem{corollary}[theorem]{Corollary}
\theoremstyle{remark}
\newtheorem{remark}[theorem]{Remark}

\newcommand{\N}{\mathbb{N}}
\newcommand{\Z}{\mathbb{Z}}
\newcommand{\sig}{\sigma}

\title{Prescribed digit sums in residue classes\\
\large and a conjecture of Murthy (OEIS A077754)}
\author{Fr\'ed\'eric Gaudin\thanks{University of Caen, France.
Email: \texttt{frederic.gaudin@unicaen.fr}}%
\thanks{2020 Mathematics Subject Classification: 11A63, 11B83.}%
\thanks{Keywords: digit sum, multiples, Murthy's conjecture,
OEIS A077754, formal verification, Lean.}}
\date{\today}

\begin{document}
\maketitle

\begin{abstract}
Let $m \ge 2$ be coprime to $10$, and put $g = \gcd(9,m)$ and $M = m/g$.
Fix a decimal length $L \ge 1$ and set $s = \lceil (M-1)/2 \rceil$. We prove that,
whenever $L \ge 3s$, \emph{every} digit sum $\sig$ in the band
$[L+3s,\,9L-3s]$ and \emph{every} residue $r \equiv \sig \pmod g$ is realized
by an $L$-digit positive integer $N$ with all digits nonzero, $\sig(N)=\sig$,
and $N \equiv r \pmod m$; this band is the feasibility window of the balanced
exchange method. The same machinery, run with denser smooth
moves and finite certificates in place of minimal moves, resolves the
conjecture of Amarnath Murthy recorded as OEIS A077754 (2002): a positive
integer $n \ge 1$ admits a multiple with at least two digits, all nonzero,
and digit sum $n$, if and only
if $n \notin \{1,2,5,11\} \cup \{10k : k \ge 1\}$. Both implications are
machine-verified in full in Lean~4 with mathlib: all
finite lemmas are kernel-checked by decide on explicit witness tables
and bitmap certificates, so the axiom base reduces to the three standard axioms
(propext, Classical.choice, Quot.sound), and an
independent self-contained Python program re-derives the main finite checks
in seconds. The OEIS sequence value
is formalized in a companion module, certifying that it is
the least admissible multiple for every non-exceptional $n \ge 1$, vanishes
on the exceptional set, and agrees with the displayed OEIS values. The proof is
elementary modulo one classical sumset theorem of Chowla (1935).
\end{abstract}

\section{Statement of results}

Throughout this paper, $\N = \{0,1,2,\dots\}$ denotes the natural numbers, and
$10\N = \{10k : k \in \N\}$ denotes the set of multiples of $10$. For a
positive integer $N \ge 1$, we write $\sig(N)$ for the sum of the decimal
digits of $N$ (with the convention $\sig(0) = 0$) and
$\operatorname{len}(N)$ for its number of decimal digits.

Entry A077754 in the On-Line Encyclopedia of Integer Sequences~\cite{oeis},
contributed by Amarnath Murthy in November 2002 for indices $n \ge 1$, reads
verbatim:
\begin{quote}
``Smallest multiple of n with two or more digits, none of them zeros, whose
digit sum equals n, or 0 if no such multiple exists. Conjecture: zero occurs
only for indices which are multiples of 10 apart from 1, 2, 5 and 11.
(a(n) = 0 only for n = 1,2,5,11 or n = 10k.)''
\end{quote}

Our first result establishes residue coverage across a wide band of digit sums
on fixed-length decimal fibres without zero digits. Murthy's conjecture is
resolved by the same balanced-exchange machinery, using denser smooth moves
and finite certificates where the minimal-move band does not reach
$\sig = n$.

\begin{theorem}[equidistribution on the digit fibre]
\label{thm:band}
Let $m \ge 2$ be coprime to $10$, and set $g = \gcd(9,m)$, $M = m/g$, and
$s = \lceil (M-1)/2 \rceil$. For every positive integer $L \ge 3s$, every integer
$\sig \in [L+3s,\,9L-3s]$, and every residue $r \equiv \sig \pmod g$, there
exists an $L$-digit positive integer $N$ with all decimal digits nonzero such
that $\sig(N)=\sig$ and $N \equiv r \pmod m$. The congruence
$r \equiv \sig \pmod g$ is necessary because $N \equiv \sig(N) \pmod g$, and
$[L+3s,\,9L-3s]$ is the feasibility window certified by the balanced exchange
method (\S\ref{sec:edges}).
\end{theorem}

\begin{theorem}[Murthy's conjecture, OEIS A077754~\cite{oeis}]
\label{cor:murthy}
A positive integer $n \ge 1$ admits a multiple $N$ with at least two decimal
digits, all nonzero, and $\sig(N) = n$, if and only if
$n \notin \{1,2,5,11\} \cup \{10k : k \ge 1\}$.
\end{theorem}

Theorem~\ref{cor:murthy} is formally verified in Lean~4 with mathlib
(v4.35.0)~\cite{lean} without \texttt{sorry} or \texttt{native\_decide},
relying only on the three standard axioms \texttt{propext},
\texttt{Classical.choice}, and \texttt{Quot.sound} (see \S\ref{sec:lean} for
the exact distinction between kernel-formalized theorems and paper-only band
arguments). Note that on $\N$, the Lean statement includes $n = 0 \in 10\N$,
which is vacuously excluded because $\operatorname{len}(0) < 2$; restricting
to $n \ge 1$ gives Theorem~\ref{cor:murthy}.

The necessity of the exceptions is elementary (\S\ref{sec:exc}). For
sufficiency, the range $1 \le n \le 200$ is settled by an exhaustive witness
table (certificate A in the appendix). For $n > 200$, we write $n = R\cdot m$
where $R$ is the $\{2,5\}$-part of $n$ and $\gcd(m,10)=1$. We concatenate a
short low-order block $B$ divisible by $R$ with a high-order prefix $H$ drawn
from a fixed-length digit fibre, realizing the target residue
$-B\cdot 10^{-d} \pmod m$ (where negative exponents denote modular inverses
modulo $m$) via balanced digit exchanges. We prove Theorem~\ref{cor:murthy}
first using dense move sets tailored to $n > 200$ (\S\ref{sec:mainproof}), and
then apply the minimal move set to establish the general band theorem
(\S\ref{sec:band}).

\begin{corollary}[quantitative length bound]
\label{cor:quant}
For every non-exceptional $n > 200$, the constructed witness $N$ satisfies
$\operatorname{len}(N) \le \lfloor n/5 \rfloor + \lfloor \log_2 n \rfloor$:
it consists of a high fibre of $L = \lfloor (n-\sig(B))/5 \rfloor$ digits
(for $m \ge 2$; in the $m = 1$ branch the fibre prefix has
$\lceil (n-\sig(B))/9 \rceil \le L$ digits)
concatenated with a low block of $d \le \log_2 n$ (or $d \le \log_5 n$) digits.
\end{corollary}

The bound in Corollary~\ref{cor:quant} is existential; the \emph{smallest}
admissible multiple -- the OEIS sequence value $a(n)$ itself -- is formalized
and certified separately for the displayed anchor values in \S\ref{sec:lean}.

\emph{Related work.} When zero digits are permitted, the existence of a
multiple of $n$ with digit sum $n$ is classical (see OEIS A002998~\cite{a002998}
and Aassila~\cite{aassila}). Related results on digit sums of multiples and
minimal Niven numbers include Fredricksen, Ionascu, Luca, and
St\u{a}nic\u{a}~\cite{fils}. None of these works enforces the all-nonzero-digit
constraint, which creates the exact exceptional set
$\{1,2,5,11\} \cup \{10k : k \ge 1\}$ in Theorem~\ref{cor:murthy}.

\section{Necessity of the exceptions}
\label{sec:exc}

We verify that no admissible multiple $N$ exists when $n \ge 1$ lies in
$\{1,2,5,11\} \cup \{10k : k \ge 1\}$.

\begin{itemize}
\item If $10 \mid n$, then every multiple $N$ of $n$ ends in the decimal digit
$0$, violating the nonzero-digit requirement.
\item If $n = 1$, every integer with at least two nonzero digits satisfies
$\sig(N) \ge 2 > 1$.
\item If $n = 2$, any multiple $N$ of $2$ with all digits nonzero must end in
an even nonzero digit (at least $2$). Having at least two nonzero digits then
forces $\sig(N) \ge 1 + 2 = 3 > 2$.
\item If $n = 5$, any multiple $N$ of $5$ with all digits nonzero must end in
$5$. Having at least two nonzero digits forces $\sig(N) \ge 1 + 5 = 6 > 5$.
\item If $n = 11$, write $S = \sig(N)$ for the digit sum of an admissible
multiple $N$ and $A$ for the alternating sum of its digits. Suppose $S = 11$
and $11 \mid N$. Then $11 \mid A$, and $A \equiv S = 11 \pmod 2$, so $A$ is
odd. Moreover, $|A| \le S = 11$. If $|A| = 11$, then all digits in either the
even or the odd decimal positions must vanish; since all digits are nonzero,
$N$ can have at most one digit, contradicting $\operatorname{len}(N) \ge 2$.
Thus $|A| < 11$, which is impossible because no odd integer strictly between
$-11$ and $11$ is divisible by $11$. \qed
\end{itemize}

These cases reflect three distinct obstructions: terminal digit divisibility
($10k$ and $5$), minimum digit-sum bounds ($1$ and $2$), and alternating digit
weight ($11$).

\section{Machinery: the fixed-count digit fibre}
\label{sec:machinery}

Let $n \ge 1$ be a positive integer not divisible by $10$. Write $n = R\cdot m$
with $\gcd(m,10)=1$ and $R \in \{1, 2^a, 5^b\}$ for integers $a,b \ge 1$; the
primes $2$ and $5$ cannot both divide $n$ because $10 \nmid n$. When $m \ge 2$,
set $g = \gcd(9,m) \in \{1,3,9\}$ and $M = m/g$.

Our witness has the form $N = B + 10^d \cdot H$, where $B$ is a $d$-digit
low-order block divisible by $R$ and $H$ is an $L$-digit high-order prefix
whose residue modulo $m$ we control by balanced exchanges. Because $B < 10^d$
has exactly $d$ nonzero digits (when $d \ge 1$), the decimal representation of
$N$ is the concatenation of $H$ and $B$, so no carries cross the interface and
$\sig(N) = \sig(H) + \sig(B)$.

\paragraph{The digit fibre.} Fix a length $L \ge 1$ and an offset sum
$K \in \{0,\dots,8L\}$. Every offset vector $u = (u_0,\dots,u_{L-1})$ with
$u_e \in \{0,\dots,8\}$ and $\sum_{e=0}^{L-1} u_e = K$ encodes an $L$-digit
integer
\[
 F(u) = \sum_{e=0}^{L-1} (1 + u_e)\,10^e .
\]
Because each digit $1 + u_e$ lies in $\{1,\dots,9\}$, no carries occur in base
$10$, all $L$ digits of $H = F(u)$ are nonzero, $\sig(H) = L + K$, and
$H \equiv L + K \pmod g$ since $10^e \equiv 1 \pmod 9$ and $g \mid 9$.

\begin{lemma}[the low block]
\label{lem:lowblock}
If $R = 1$, set $d = 0$ and $B = 0$. If $R = 2^a$ with $a \ge 1$, there exists
an $a$-digit integer $B$ with all digits in $\{1,2\}$ such that $2^a \mid B$
and $\sig(B) \le 2a$. If $R = 5^b$ with $b \ge 1$, there exists a $b$-digit
integer $B$ with all digits in $\{1,2,3,4,5\}$ such that $5^b \mid B$ and
$\sig(B) \le 5b$.
\end{lemma}

\begin{proof}
For $R = 2^a$, proceed by induction on $a$. For $a = 1$, take $B = 2$. Given a
$j$-digit multiple $B_j$ of $2^j$ with digits in $\{1,2\}$, we have
$10^j = 2^j \cdot 5^j \equiv 2^j \pmod{2^{j+1}}$, so one of $B_j + 1\cdot 10^j$
or $B_j + 2\cdot 10^j$ is divisible by $2^{j+1}$. Similarly, for $R = 5^b$,
take $B_1 = 5$; given $5^j \mid B_j$, write $10^j = 5^j \cdot 2^j$. Since
$\gcd(2^j,5)=1$, as $c$ ranges over $\{1,2,3,4,5\}$ the values
$B_j + c\cdot 10^j$ cover all five multiples of $5^j$ modulo $5^{j+1}$, so one
is divisible by $5^{j+1}$. The digit-sum bounds follow immediately.
\end{proof}

\paragraph{Balanced exchanges.} On a block of three consecutive positions
$\{3i,3i+1,3i+2\}$ inside $\{0,\dots,L-1\}$, consider a triple
$(d_0,d_1,d_2) \in \{-4,\dots,4\}^3$ such that $d_0 = -d_1 - d_2$ and
$u_{3i+r} + d_r \in \{0,\dots,8\}$ for $r \in \{0,1,2\}$. Adding $(d_0,d_1,d_2)$
to $(u_{3i},u_{3i+1},u_{3i+2})$ preserves the total sum $\sum_e u_e = K$, keeps
all digits in $\{1,\dots,9\}$, and shifts the value $F(u)$ by
\[
 \Delta = (d_0 + 10 d_1 + 100 d_2)\cdot 10^{3i}
 = 9\cdot 10^{3i}\cdot(d_1 + 11 d_2).
\]
When $u_{3i} = u_{3i+1} = u_{3i+2} = 4$, any triple with $d_1, d_2, -d_1-d_2 \in
\{-4,\dots,4\}$ is admissible; the resulting integer grid
\[
 X = \{\, d_1 + 11 d_2 : d_1, d_2 \in \{-4,\dots,4\},\ |d_1+d_2| \le 4 \,\}
\]
contains $61$ elements ($0$ together with $60$ nonzero integers, all bounded in
absolute value by $44$). When we only assume $u_{3i+r} \in \{1,\dots,7\}$ for
$r \in \{0,1,2\}$, restricting to $d_0,d_1,d_2 \in \{-1,0,1\}$ with
$d_0 = -d_1 - d_2$ yields $d_1 + 11 d_2 \in \{0, \pm 1, \pm 10, \pm 11\}$; in
particular, the \emph{minimal moves} $x \in \{\pm 1, \pm 10\}$ are always
available on any triple of positions with offsets in $\{1,\dots,7\}$.

\paragraph{Frozen poles.} An offset $u_e = 0$ (digit $1$) admits only $d_r \ge 0$,
and an offset $u_e = 8$ (digit $9$) admits only $d_r \le 0$. If all three
positions of a block sit at $u = 0$ (or all three sit at $u = 8$), the balance
condition $d_0 + d_1 + d_2 = 0$ forces $d_0 = d_1 = d_2 = 0$. Requiring every
position in each active exchange block to satisfy $u_e \in \{1,\dots,7\}$ avoids
both frozen poles uniformly.

\section{Stage sets and the Chowla step}
\label{sec:stages}

Let $m \ge 2$ be coprime to $10$, $g = \gcd(9,m)$, and $M = m/g$. Because
$\gcd(m,10) = 1$, $M$ is coprime to $10$ (in particular, $M$ is odd). Moreover,
$\gcd(9/g, M) = 1$: this is trivial for $g \in \{1,9\}$, and if $g = 3$ then
$9 \nmid m$, so $3 \nmid (m/3) = M$.

If $M = 1$ (which occurs when $m \in \{3,9\}$), the ring $\Z/M\Z$ is trivial:
every integer is congruent to $0 \pmod M$, so $s = 0$ stages already cover
$\Z/M\Z$ once the congruence modulo $g = m$ holds.

Assume henceforth that $M \ge 2$ (and hence $M \ge 3$, since $\gcd(M,10)=1$).
For a subset $U \subseteq X$ of grid values that are admissible on block $i$,
we define the \emph{stage set} modulo $M$ by
\[
 B^*_i(U) = \{0\} \cup \{\, (9/g)\cdot 10^{3i}\cdot x \bmod M :
 x \in U,\ \gcd(x,M) = 1 \,\} \subseteq \Z/M\Z .
\]
Because $\gcd(9/g,M) = \gcd(10^{3i},M) = 1$, multiplication by $(9/g)\cdot 10^{3i}$
is a unit automorphism of $\Z/M\Z$: every nonzero element of $B^*_i(U)$ is a
unit in $\Z/M\Z$, and $|B^*_i(U)| = |B^*_0(U)|$ is independent of $i$. Distinct
blocks $i$ may produce the same subset of $\Z/M\Z$ (for instance when $10^3
\equiv 1 \pmod M$, as for $M = 333$), which causes no difficulty because
different blocks act on disjoint triples of decimal positions.

\begin{theorem}[Chowla, 1935~\cite{chowla}]
\label{thm:chowla}
Let $k \ge 2$ be an integer, and let $A, B \subseteq \Z/k\Z$ be nonempty subsets.
If $0 \in B$ and every nonzero element of $B$ is coprime to $k$, then
$|A+B| \ge \min(k,\,|A|+|B|-1)$.
\end{theorem}

Applying Theorem~\ref{thm:chowla} iteratively across $s$ distinct position
blocks $i_1, \dots, i_s$ (which requires $3s$ positions in total), starting
from $\{0\} \subseteq \Z/M\Z$, yields
\[
 |B^*_{i_1}(U) + \cdots + B^*_{i_s}(U)| \ge \min(M,\, s\cdot w + 1),
 \qquad w = |B^*_0(U)| - 1 .
\]
Hence $s = \lceil (M-1)/w \rceil$ stages cover all of $\Z/M\Z$ whenever $w \ge 1$.
We apply this in two regimes (still assuming $M \ge 3$):

\begin{itemize}
\item \emph{Full smooth moves} (base offsets $u_e = 4$ on the active blocks,
admitting all of $X$): when $g \in \{3,9\}$, the $\{2,5\}$-smooth elements of
$X \setminus \{0\}$ are
\[
 U_{2,5} = \{\pm 1, \pm 2, \pm 4, \pm 8, \pm 10, \pm 20, \pm 32, \pm 40\},
\]
a set of $16$ integers automatically coprime to $M$ since $\gcd(M,10)=1$. When
$g = 1$, we also have $3 \nmid M$, so we may adjoin the $\{2,3,5\}$-smooth
elements $\pm\{3, 9, 12, 18, 24, 30\} \subseteq X$, obtaining a $28$-element
set $U_{2,3,5} \subseteq X$ of integers coprime to $M$ (certificate C). All
elements of $U_{2,5}$ and $U_{2,3,5}$ lie in $[-40,40]$, so their pairwise
differences are nonzero and bounded in absolute value by $80$. Therefore,
\emph{whenever $M > 80$}, reduction modulo $M$ is injective on $U_{2,5}$ and on
$U_{2,3,5}$ (and none of these units can reduce to $0 \bmod M$), giving
$w = |B^*_0(U_{2,5})| - 1 = 16$ for $g \in \{3,9\}$ and
$w = |B^*_0(U_{2,3,5})| - 1 = 28$ for $g = 1$. Consequently, for $M > 80$, at
most $s = \lceil (M-1)/w \rceil$ stages cover $\Z/M\Z$. (The Lean
formalization uses the slightly more conservative uniform step count
$s = \lfloor m/27 \rfloor + 1$ -- that is, $\lfloor M/27 \rfloor + 1$,
$\lfloor M/9 \rfloor + 1$, or $\lfloor M/3 \rfloor + 1$ according as
$g = 1$, $3$, or $9$ -- which satisfies the
same $3s + 4 \le L$ inequality for all $M > 184$.)
\item \emph{Minimal moves} (offsets $u_e \in \{1,\dots,7\}$ on the active
blocks): take $U_{\min} = \{\pm 1, \pm 10\}$. Since $\gcd(M,10)=1$, all four
elements of $U_{\min}$ are units modulo $M$. Because $M \ge 3$ is odd and
$c = (9/g)\cdot 10^{3i}$ is a unit modulo $M$, the residues $+c$ and $-c$ are
distinct modulo $M$ (otherwise $M \mid 2c$, impossible since $\gcd(2c,M)=1$).
Thus $B^*_i(U_{\min}) \supseteq \{0, +c, -c\}$ satisfies $|B^*_i(U_{\min})| \ge 3$,
so $w \ge 2$ for every $M \ge 3$, and $s = \lceil (M-1)/2 \rceil$ stages cover
$\Z/M\Z$. (In fact, adding $\{0, +c, -c\}$ to any nonempty proper subset
$A \subsetneq \Z/M\Z$ strictly increases $A$ by at least $2$ elements up to
size $M$ by considering orbits of $+c$, giving Theorem~\ref{thm:chowla}
directly in this special case.)
\end{itemize}

\section{Proof of Theorem~\ref{cor:murthy}}
\label{sec:mainproof}

Necessity was proved in \S\ref{sec:exc}, and every non-exceptional $n$ in the
range $1 \le n \le 200$ is resolved by an explicit witness in certificate A.
Assume now that $n > 200$ and $10 \nmid n$. Write $n = R\cdot m$ with
$\gcd(m,10)=1$ and $R \in \{1, 2^a, 5^b\}$, and let $B$ and $d$ be given by
Lemma~\ref{lem:lowblock}.

\paragraph{Case $m = 1$.} Here $n = R > 200$. If $n = 2^a$, then $a \ge 8$ and
$\sig(B) \le 2a < 2^a = n$. If $n = 5^b$, then $b \ge 4$ and
$\sig(B) \le 5b < 5^b = n$. Thus $S := n - \sig(B) \ge 1$ in both cases. Write
$S = 9q + r$ with $q \ge 0$ and $r \in \{1,\dots,9\}$, and let $H$ be the
integer consisting of the digit $r$
followed by $q$ digits equal to $9$. All digits of $H$ are nonzero,
$\sig(H) = S$, and $N = B + 10^d \cdot H$ has all digits nonzero,
$\operatorname{len}(N) \ge d + 1 \ge 5 \ge 2$, $\sig(N) = \sig(B) + S = n$,
and $R \mid N$ since $R \mid B$ and $R \mid 10^d$.

\paragraph{Case $m \ge 2$: target and coset compatibility.} Since $\gcd(m,10)=1$,
$10^d$ is invertible modulo $m$; set the target residue
$t := -B \cdot 10^{-d} \bmod m$ (so $t = 0$ when $R = 1$). We choose the fibre
parameters
\[
 L = \Big\lfloor \frac{n - \sig(B)}{5} \Big\rfloor, \qquad
 K = n - \sig(B) - L = 4L + r_0, \quad r_0 \in \{0,1,2,3,4\},
\]
provided $n - \sig(B) \ge 5$ (verified below). Because $B \equiv \sig(B) \pmod 9$
and $10^d \equiv 1 \pmod 9$, we have $t \equiv -\sig(B) \pmod g$. On the fibre,
every encoded value $F(u)$ satisfies
$F(u) \equiv L + K = n - \sig(B) \equiv -\sig(B) \equiv t \pmod g$, because
$g \mid m$ and $m \mid n$. Thus $(t - F(u^0))/g$ is an integer for any base
offset vector $u^0$ of sum $K$.

To place $s$ disjoint exchange blocks at base offset $u_e = 4$ and distribute
the remainder $r_0 \le 4$ by raising $r_0$ separate positions from $4$ to $5$,
it suffices that $L \ge 3s + 4$ (or $L \ge 3s + r_0$).

\paragraph{Bounding $\sig(B)$ and $n - \sig(B)$.} If $R = 1$, then $\sig(B) = 0$
and $n - \sig(B) = m = gM$. Suppose $R \ge 2$. If $R = 2^a$ ($a \ge 1$), then
$2^a - 2a \ge 0$ for all $a \ge 1$; if $R = 5^b$ ($b \ge 1$), then
$5^b - 5b \ge 0$ for all $b \ge 1$. Moreover, since $m \ge 3$ (as $\gcd(m,10)=1$
and $m \ge 2$), we have
\[
 n - \sig(B) = R m - \sig(B) = m + (R-1)m - \sig(B) \ge gM + 3(R-1) - \sig(B).
\]
For $R = 2^a$ ($a \ge 1$), $3(2^a - 1) - 2a \ge 3\cdot 1 - 2 = 1 > 0$; for
$R = 5^b$ ($b \ge 1$), $3(5^b - 1) - 5b \ge 3\cdot 4 - 5 = 7 > 0$. Therefore,
in \emph{all} cases ($R = 1$ or $R \ge 2$), we have the clean lower bound
\[
 n - \sig(B) \ge m = gM .
\]
Similarly, because $n > 200$, one checks directly that $n - \sig(B) \ge 188$:
for $R = 1$ we have $n - \sig(B) = n \ge 201$; for $R = 2^a$ with $n = 2^a m > 200$
and $m \ge 3$, we have $2^a m - 2a \ge 202 - 14 = 188$ (and if $a \ge 8$,
$2^a m - 2a \ge 3\cdot 256 - 16 = 752$); for $R = 5^b$ with $n = 5^b m > 200$
and $m \ge 3$, the smallest value of $5^b m - 5b$ for $b \ge 2$ occurs at
$b = 2, m = 9$
($n = 225$), giving $225 - 10 = 215$ (the next case $b = 3, m = 3$ gives
$375 - 15 = 360$, and $b \ge 4$ gives $\ge 1855$), while $b = 1$ has $m \ge 41$
and $5m - 5 \ge 200$. Hence
$n - \sig(B) \ge 188$ for every $n > 200$ with $m \ge 2$, which forces
\[
 L = \Big\lfloor \frac{n - \sig(B)}{5} \Big\rfloor \ge \Big\lfloor \frac{188}{5} \Big\rfloor = 37 .
\]

\paragraph{Verifying the stage budget $3s + 4 \le L$.} We split on the size of $M$.

\begin{itemize}
\item \emph{Analytic regime $M > 184$.} Since $M > 184 > 80$, the smooth-unit
injectivity holds (\S\ref{sec:stages}), giving $w = 28$ when $g = 1$ and
$w = 16$ when $g \in \{3,9\}$, with $s = \lceil (M-1)/w \rceil$. Using
$L = \lfloor (n - \sig(B))/5 \rfloor \ge \lfloor gM/5 \rfloor \ge (gM - 4)/5$:
\begin{itemize}
\item For $g = 1$ ($w = 28$, $M \ge 185$): we need
$3\lceil (M-1)/28 \rceil + 4 \le \lfloor M/5 \rfloor$. Since
$\lceil (M-1)/28 \rceil \le (M + 26)/28$, we have
\[
 3\Big\lceil \frac{M-1}{28} \Big\rceil + 4 \le \frac{3M + 190}{28} \le \frac{M - 4}{5}
\]
for all $M \ge 184$ (as $15M + 950 \le 28M - 112 \iff 13M \ge 1062 \iff M \ge 82$).
\item For $g = 3$ ($w = 16$, $M \ge 185$): we have
\[
 3\Big\lceil \frac{M-1}{16} \Big\rceil + 4 \le \frac{3M + 106}{16} \le \frac{3M - 4}{5}
\]
for all $M \ge 185$ (as $15M + 530 \le 48M - 64 \iff 33M \ge 594 \iff M \ge 18$).
\item For $g = 9$ ($w = 16$, $M \ge 185$): we have
$3\lceil (M-1)/16 \rceil + 4 \le (3M - 4)/5 < (9M - 4)/5 \le L$.
\end{itemize}
Thus $3s + 4 \le L$ holds for all $M > 184$.

\item \emph{Certified finite regime $M \le 184$.} Here $m = gM \le 9 \cdot 184 = 1656$.
If $M = 1$ ($m \in \{3,9\}$), $s = 0$ stages suffice and $4 \le 37 \le L$. For
$2 \le M \le 184$ with $m \notin \{11,19,23\}$, certificate B verifies directly
by bit-parallel sumset evaluation that the smooth moves $U_{2,3,5}$ ($28$ values,
used when $g = 1$) or $U_{2,5}$ ($16$ values, used when $g \in \{3,9\}$) cover
$\Z/m\Z$ (within
the $\bmod\, g$ coset) in $k \le \lfloor m/27 \rfloor + 1$ stages. Using
$n - \sig(B) \ge m$ and $L \ge 37$: if $m \le 296$, then $k \le 11$ and
$3k + 4 \le 37 \le L$; if $m \ge 297$, then
\[
 3k + 4 \le 3\Big(\frac{m}{27} + 1\Big) + 4 = \frac{m}{9} + 7 \le \frac{m - 4}{5} \le L
\]
where the last step uses the stronger available bound $L \ge \lfloor m/5 \rfloor$
since $5m + 315 \le 9m - 36 \iff 4m \ge 351 \iff m \ge 88$. Finally, for the
three small moduli $m \in \{11,19,23\}$, certificate D1 verifies that $k = 2$
stages cover $\Z/m\Z$ (requiring $3k + 4 = 10 \le 37 \le L$). (In the Lean
formalization, the standalone lemma for $m \ge 2$ is stated more generally for
all $L \ge \lfloor m/5 \rfloor$ rather than only $L \ge 37$; for
$m \in \{11,19,23\}$ at the small lengths $L \in \{7,8,9\}$, where $3k+4=10 > L$,
certificate D2 supplies an explicit table of $795$ witnesses.)
\end{itemize}

\paragraph{Assembly.} Let $u^0$ be the base offset vector with $u^0_e = 4$ on
the first $3s$ positions (forming $s$ blocks), $u^0_e = 5$ on $r_0$ of the
remaining positions, and $u^0_e = 4$ elsewhere, and set $H_0 = F(u^0)$. By
stage coverage, there exist $s$ admissible grid values $x_0,\dots,x_{s-1} \in X$
such that
\[
 \sum_{i=0}^{s-1} (9/g)\,10^{3i}\,x_i \equiv \frac{t - H_0}{g} \pmod M .
\]
Lifting each $x_i = d_{i,1} + 11 d_{i,2}$ to a balanced triple on block $i$
produces an offset vector $u \in \{0,\dots,8\}^L$ with $\sum_e u_e = 4L + r_0 = K$
and $H := F(u) = H_0 + g \sum_{i=0}^{s-1} (9/g)\,10^{3i}\,x_i \equiv t \pmod m$.
Then $N = B + 10^d \cdot H$ has all digits nonzero, $\operatorname{len}(N) = d + L \ge 37 \ge 2$,
and $\sig(N) = \sig(B) + L + K = n$. Moreover, $N \equiv 0 \pmod R$ (since
$R \mid B$ and $R \mid 10^d$) and $N \equiv B + 10^d t \equiv 0 \pmod m$; since
$\gcd(R,m)=1$, we conclude $n = Rm \mid N$. \qed

\section{Proof of Theorem~\ref{thm:band}}
\label{sec:band}

Let $m \ge 2$ be coprime to $10$, $g = \gcd(9,m)$, $M = m/g$, and
$s = \lceil (M-1)/2 \rceil$. Fix a positive integer $L \ge 3s$, an
integer $\sig \in [L+3s,\,9L-3s]$, and a residue
$r \equiv \sig \pmod g$. Set $K = \sig - L$, so $K \in [3s,\,8L-3s]$.

If $M = 1$ (that is, $m \in \{3,9\}$ and $g = m$), then $s = 0$ and
$K \in [0,\,8L]$. Any offset vector $u \in \{0,\dots,8\}^L$ with $\sum_e u_e = K$
(obtained greedily by filling offsets up to $8$) gives an $L$-digit nonzero-digit
integer $N = F(u)$ with $\sig(N) = L + K = \sig$ and
$N \equiv \sig \equiv r \pmod m$.

Assume now that $M \ge 3$, so $s = \lceil (M-1)/2 \rceil \ge 1$ and $L \ge 3s \ge 3$.

\paragraph{(i) Offset distribution.} Initialize $u^0_e = 1$ on the $3s$ active
positions $\{0,\dots,3s-1\}$ (which uses $3s$ units of sum) and $u^0_e = 0$ on
the remaining $L - 3s$ positions $\{3s,\dots,L-1\}$. We must distribute the
remaining $K - 3s \ge 0$ units across the $L$ positions while keeping
$u^0_e \in \{1,\dots,7\}$ for $e \in \{0,\dots,3s-1\}$ and
$u^0_e \in \{0,\dots,8\}$ for $e \in \{3s,\dots,L-1\}$. Each of the $3s$ active
positions can absorb up to $6$ additional units (from $1$ up to $7$), and each
of the $L - 3s$ free positions can absorb up to $8$ additional units (from $0$
up to $8$). The total remaining capacity is
\[
 6(3s) + 8(L - 3s) = 8L - 6s .
\]
Because $K \le 8L - 3s$, we have $K - 3s \le 8L - 6s$, so a greedy allocation
produces a base vector $u^0$ with $\sum_{e=0}^{L-1} u^0_e = K$ and
$u^0_e \in \{1,\dots,7\}$ for all $e \in \{0,\dots,3s-1\}$.

\paragraph{(ii) Residue coverage.} Let $H_0 = F(u^0)$. Since $H_0 \equiv L + K = \sig \equiv r \pmod g$,
the integer $(r - H_0)/g$ is well-defined modulo $M$. Because $u^0_e \in \{1,\dots,7\}$
on every position of the $s$ disjoint blocks $\{3i,3i+1,3i+2\}$ ($0 \le i < s$),
the minimal moves $U_{\min} = \{\pm 1, \pm 10\}$ are admissible on each block.
As shown in \S\ref{sec:stages}, for $M \ge 3$ each stage set $B^*_i(U_{\min})$
has cardinality at least $3$ (so $w \ge 2$) and every nonzero element is a unit
modulo $M$. By Theorem~\ref{thm:chowla}, the $s$-fold sumset
$B^*_0(U_{\min}) + \cdots + B^*_{s-1}(U_{\min})$ has size at least
$\min(M,\, 2s + 1) = M$, covering $\Z/M\Z$. Selecting moves $x_i \in U_{\min} \cup \{0\}$
for $0 \le i < s$ with $\sum_{i=0}^{s-1} (9/g)\,10^{3i}\,x_i \equiv (r - H_0)/g \pmod M$
and applying the corresponding balanced triples in $\{-1,0,1\}^3$ yields an
offset vector $u \in \{0,\dots,8\}^L$ of sum $K$ such that $N = F(u)$ has $L$
nonzero digits, $\sig(N) = \sig$, and $N \equiv r \pmod m$. \qed

\begin{remark}
One cannot uniformly replace $s = \lceil (M-1)/2 \rceil$ in the minimal-move
argument by $\lceil (M-1)/4 \rceil$: when $M \in \{9,11\}$, the congruence
$10 \equiv \pm 1 \pmod M$ collapses $U_{\min} \bmod M$ to $\{\pm 1\}$, leaving
$|B^*_i(U_{\min})| = 3$ ($w = 2$) so that $4$ stages (for $M = 9$) and $5$
stages (for $M = 11$) are genuinely required by minimal moves. An auxiliary
machine check over all $m \le 200$ coprime to $10$ at minimal non-vacuous
lengths confirms Theorem~\ref{thm:band} on every residue and digit sum in the
stated band; this computation is not part of the certified development.
\end{remark}

\section{Method window and edge behavior}
\label{sec:edges}

We clarify the role of the endpoints of the interval $[L+3s,\,9L-3s]$ in
Theorem~\ref{thm:band} (where $L \ge 1$ is assumed). To operate $s$ independent minimal-move exchange blocks
simultaneously without risking frozen poles, every one of the $3s$ active
positions must lie in $\{1,\dots,7\}$. Having at least $3s$ positions with
$u_e \ge 1$ forces $K = \sum_e u_e \ge 3s$, i.e.\ $\sig \ge L + 3s$;
symmetrically, having at least $3s$ positions with $u_e \le 7$ forces
$8L - K = \sum_e (8 - u_e) \ge 3s$, i.e.\ $\sig \le 9L - 3s$. In this precise
sense, $[L+3s,\,9L-3s]$ is the exact feasibility window of our $s$-stage
non-frozen block construction.

We do \emph{not} claim that residues outside $[L+3s,\,9L-3s]$ are
mathematically unrealizable: individual moduli $m$ and sums $\sig$ near the
edges can admit witnesses via denser moves from $X$ or asymmetric position
choices. Only at the extreme endpoints $\sig = L$ (the repunit $11\cdots 1$)
and $\sig = 9L$ (the number $99\cdots 9$) does the fibre collapse to a single
integer with no residue freedom at all. Inside $[L+3s,\,9L-3s]$, disjointness
of the $s$ position triples prevents any interference between stages regardless
of the multiplicative order of $10$ modulo $M$.

\section{Formal verification}
\label{sec:lean}

It is important to distinguish what is kernel-verified in Lean~4 from what
belongs to the mathematical presentation of this paper:
\begin{itemize}
\item \emph{Kernel-verified in Lean~4 (v4.35.0, mathlib v4.35.0)~\cite{lean}:}
the complete characterization of Murthy's conjecture (Theorem~\ref{cor:murthy}),
all five necessity lemmas, the fixed-sum smooth-exchange machinery feeding the
existence proof for $n > 200$, all six finite certificates (A, B, C, D1, D2,
and E), and the OEIS sequence-value specification (\texttt{IsLeast} for
non-exceptional indices, vanishing on exceptional indices, and exact values on
displayed OEIS anchors).
\item \emph{Proved in this paper only:} the general two-parameter band theorem
for arbitrary $\sig \in [L+3s,\,9L-3s]$ using minimal moves
(Theorem~\ref{thm:band}, \S\ref{sec:band}) and the quantitative summary
(Corollary~\ref{cor:quant}).
\end{itemize}

In the main file \texttt{A077754.lean}, the headline theorem is formalized over
$n \in \N$ as the equivalence between the existence of $N \in \N$ satisfying
$n \mid N$, $\mathrm{allNonzero}(N)$, $\sig_{10}(N) = n$, and
$(\mathrm{Nat.digits}\ 10\ N).\mathrm{length} \ge 2$, and the condition
$n \notin \{1,2,5,11\} \cup \{m \in \N : 10 \mid m\}$. Table~\ref{tab:lean}
records the correspondence between results in the text and declarations in the
Lean artifact.

\begin{table}[ht]
\centering
\small
{\footnotesize\setlength{\tabcolsep}{4pt}
\begin{tabular}{llp{6.2cm}}
\hline
Paper item & Lean declaration & Role in formalization \\
\hline
Theorem~\ref{cor:murthy} & \texttt{A077754} & Full iff characterization of Murthy's conjecture \\
\S\ref{sec:exc} ($n=1,2,5,10k,11$) & \texttt{exc\_1}, \texttt{exc\_2}, \texttt{exc\_5}, \texttt{exc\_10}, \texttt{exc\_11} & Necessity of the exceptional set \\
\S\ref{sec:mainproof} ($n \le 200$, Cert.\ A) & \texttt{existence\_small} & Exhaustive witness table for $n \le 200$ \\
\S\ref{sec:mainproof} (inductive / low block) & \texttt{existence\_step} & Reduction $n = R m$ via low-order block \\
\S\ref{sec:mainproof} (coprime factor $m$) & \texttt{existence\_Rm} & High-fibre existence modulo $m$ ($\gcd(m,10)=1$) \\
\S\ref{sec:machinery} (exchange grid $X$) & \texttt{Xgrid\_card} & Grid structure and cardinality check \\
\S\ref{sec:stages} (Cert.\ C, smooth units) & \texttt{Xsmooth25\_sub}, \texttt{Xsmooth235\_sub} & Inclusion and counts of smooth grid units \\
\S\ref{sec:mainproof} (Cert.\ B, $M \le 184$) & \texttt{stage\_cover\_small\_aux'} & Bit-parallel stage coverage for $m \le 1656$ \\
\S\ref{sec:mainproof} (Cert.\ D2, corners) & \texttt{corner\_witness} & $795$ explicit witnesses for $m \in \{11,19,23\}$, $L \in \{7,8,9\}$ \\
\S\ref{sec:lean} (OEIS sequence $a(n)$) & \texttt{a077754\_val}, \texttt{a077754\_val\_spec} & Definition of $a(n)$ and \texttt{IsLeast} specification \\
\S\ref{sec:lean} (Cert.\ E, zero anchors) & \texttt{a077754\_val\_zero\_*} & Certified zeros at displayed exceptional indices \\
\hline
\end{tabular}}
\caption{Correspondence between mathematical statements and Lean~4 declarations.}
\label{tab:lean}
\end{table}

In the companion module \texttt{A077754Val.lean}, the OEIS sequence itself is
defined via \texttt{a077754\_val}: \texttt{a077754\_val\_spec} certifies that
for every non-exceptional $n$, \texttt{a077754\_val}\,$n$ is the least element
(\texttt{IsLeast}) of the set of admissible multiples, while vanishing exactly
on the exceptional set. The displayed OEIS terms are kernel-certified value
anchors: $a(3)=12$, $a(4)=112$, $a(6)=24$, $a(7)=133$, $a(8)=152$, $a(9)=18$,
and $a(12)=48$ (each verified together with minimality by scanning all positive
multiples $n\cdot j < a(n)$), alongside $a(n)=0$ at
$n \in \{1,2,5,10,11,20,30,40\}$.

Running \texttt{lake build} compiles $9052$ jobs with $0$ errors and no
\texttt{sorry} across four Lean modules (approximately $2600$ lines in
\texttt{A077754.lean}, $267$ lines in \texttt{A077754Cert.lean}, $805$ lines in
\texttt{CornerTable.lean}, and $319$ lines in \texttt{A077754Val.lean}). Every
finite certificate is evaluated inside the Lean kernel via \texttt{decide},
avoiding \texttt{native\_decide} so that the axiom report
for \texttt{A077754} lists only
\texttt{propext}, \texttt{Classical.choice}, and \texttt{Quot.sound}. For the
sumset coverage certificate B, the kernel evaluates a bit-parallel residue
bitmap in the dedicated module \texttt{A077754Cert.lean}, and a bridge lemma in
\texttt{A077754.lean} transports the bitmap verdict to \texttt{Finset}
coverage.

\paragraph{Artifact.} This paper, the Lean formalization
(modules \texttt{A077754}, \texttt{A077754Cert},
\texttt{CornerTable}, \texttt{A077754Val}), the independent
verifier script, and the pinned build files are archived
together at \url{https://doi.org/10.5281/zenodo.23879087}.

\section{Remarks}

\begin{itemize}
\item \emph{Alternative asymptotic tools.} Let $Z_j = (10^j - 1)/9$ denote the
repunit of $j$ ones. A two-block expression of the form
$Z_a \cdot 10^k + (10^k - Z_r)$ (with $1 \le r \le k$, which consists of $a$
ones, followed by $k-r$ nines, $r-1$ eights, and a final nine, hence all digits
nonzero and controllable digit sum) can be combined with Konyagin's results on
arithmetic properties of integers with missing digits~\cite{konyagin} to treat
sufficiently large $n$. The present proof replaces heavy exponential-sum inputs
by Chowla's elementary theorem together with explicit finite certificates that
close the threshold down to $n = 1$.
\item \emph{Other bases.} In a general base $b \ge 4$, the same fixed-length
fibre and balanced triple exchange mechanism applies directly (with the prime
factors of $b$ playing the role of $\{2,5\}$, and digits in $\{1,\dots,b-1\}$
leaving $b - 3 \ge 1$ non-frozen interior values $\{2,\dots,b-2\}$ (at least
two for $b \ge 5$)). Base $b = 2$
degenerates completely since $1$ is the unique nonzero digit, and base $b = 3$
has no non-frozen interior digit either (its nonzero digits $\{1,2\}$ are both
frozen poles), so base $3$ requires a redesigned multi-position exchange gadget
rather than single position-triple offsets.
\item \emph{Role of the dense exchange grid.} Balanced digit exchanges and
Chowla's 1935 sumset inequality are classical tools, but using minimal moves
alone ($w = 2$) would require $L \approx 3M/2$ digits, exceeding the available
linear digit budget $L \approx n/5$ when $g = 1$ and $n = M$. Enlarging the
single-block grid $X$ to $60$ nonzero values yields $28$ smooth units when
$g = 1$ and $16$ when $g \in \{3,9\}$, reducing the stage slope to $3/28 < 1/5$
(respectively $3/16 < 3/5$) and pushing the analytic threshold down to
$M \le 184$, well within reach of kernel-checked bitmap evaluation.
\end{itemize}

\appendix
\section{Reproducible machine certificates}

Six finite certificates are used in the formalization; all six are
kernel-checked by \texttt{decide} inside the Lean artifact (\S\ref{sec:lean}),
and certificates A, B, and C are independently re-derived in seconds by the
standalone script \texttt{A077754-ANNEXE.py} (Python standard library, around
$90$ lines):

\begin{itemize}
\item \emph{Certificate A (small $n$ witnesses):} an explicit admissible
multiple $N$ for each of the $176$ non-exceptional positive integers
$n \le 200$, found by breadth-first search over states
$(\text{residue} \bmod n,\, \text{digit sum},\, \text{length})$.
\item \emph{Certificate B (small-$M$ stage coverage):} in
\texttt{A077754Cert.lean}, bit-parallel sumset evaluation verifies that for
every $m \le 1656$ coprime to $10$ with $2 \le M \le 184$ and
$m \notin \{11,19,23\}$, the smooth move set $U_{2,3,5}$ ($28$ values, used when
$g = 1$) or $U_{2,5}$ ($16$ values, used when $g \in \{3,9\}$) covers all
residues in the
appropriate $\bmod\, g$ coset of $\Z/m\Z$ within $k \le \lfloor m/27 \rfloor + 1$
stages. (The Python script also runs a uniform check showing that every
$M \le 600$ coprime to $10$ is covered in at most $6$ stages, with worst case
$M = 333$ where $10^3 \equiv 1 \pmod{333}$.)
\item \emph{Certificate C (smooth grid units):} verification on the $61$-element
grid $X$ that $X$ contains the $16$ $\{2,5\}$-smooth nonzero values
$\pm\{1,2,4,8,10,20,32,40\}$ (used for $g \in \{3,9\}$) together with the $12$
additional $\{2,3,5\}$-smooth values $\pm\{3,9,12,18,24,30\}$, giving $28$
smooth units for $g = 1$.
\item \emph{Certificate D1 (two-stage coverage for $\{11,19,23\}$):} a finite
stage-coverage check (not a witness table) confirming that $k = 2$ stages of
smooth moves from $U_{2,3,5}$ (which applies since $g = \gcd(9,m) = 1$ for all
three moduli) cover $\Z/m\Z$ for $m \in \{11,19,23\}$ (in fact $k = 1$ already
covers $m = 19$ and $m = 23$).
\item \emph{Certificate D2 (corner witness table):} in \texttt{CornerTable.lean}
and \texttt{A077754.lean}, an explicit table of $795$ digit witnesses covering
all residues modulo $m$ for $m \in \{11,19,23\}$, $L \in \{7,8,9\}$, and
$r_0 \in \{0,1,2,3,4\}$. As noted in \S\ref{sec:mainproof}, $n > 200$ forces
$L \ge 37$ in the main theorem, so D2 is used solely to close the standalone
helper lemma at its stated lower bound $L \ge \lfloor m/5 \rfloor$.
\item \emph{Certificate E (OEIS value anchors):} in \texttt{A077754Val.lean},
kernel verification of the exact sequence values $a(3)=12$, $a(4)=112$,
$a(6)=24$, $a(7)=133$, $a(8)=152$, $a(9)=18$, and $a(12)=48$ (each paired with
a kernel scan over $n \cdot j < a(n)$ proving minimality, distinguishing e.g.\
the minimal value $a(7) = 133$ from non-minimal existential witnesses), together
with $a(n) = 0$ for $n \in \{1,2,5,10,11,20,30,40\}$.
\end{itemize}

\section*{Use of AI tools}
Generative AI tools assisted in blind OEIS transcription,
adversarial auditing, proof sketching, and text review,
within a verification methodology designed and supervised
by the author. The characterization of OEIS A077754
(Theorem~\ref{cor:murthy}) and its sequence-value
specification are verified from standard axioms only in
the Lean~4 kernel, with the Lean source and the
independent Python certificate script archived in the
companion record. The author selected the problems and
proof strategies, arbitrated all discrepancies, verified
OEIS fidelity, and takes full responsibility for this
work.

\section*{Funding}
No funding was received for this work.

\section*{Competing interests}
The author declares none.

\begin{thebibliography}{9}
\bibitem{oeis} A.~Murthy, OEIS A077754 (conjecture,
 Nov.~2002), in \emph{The On-Line Encyclopedia of Integer
 Sequences}, N.~J.~A.~Sloane, ed., \url{https://oeis.org/A077754}.
\bibitem{a002998} OEIS A002998, in \emph{The On-Line
 Encyclopedia of Integer Sequences}, N.~J.~A.~Sloane,
 ed., \url{https://oeis.org/A002998}.
\bibitem{chowla} I.~Chowla, \emph{A theorem on the addition
 of residue classes: Application to the number $\Gamma(k)$ in
 Waring's problem}, Proc.\ Indian Acad.\ Sci.\ Sect.~A
 \textbf{2} (1935), 242--243.
\bibitem{konyagin} S.~Konyagin, \emph{Arithmetic properties
 of integers with missing digits}, Period.\ Math.\ Hungar.\
 \textbf{42} (2001), 145--162.
\bibitem{fils} H.~Fredricksen, E.~J.~Ionascu, F.~Luca, and
 P.~St\u{a}nic\u{a}, \emph{Minimal Niven numbers}, Acta Arith.\
 \textbf{132} (2008), 135--159; arXiv:0803.0477.
\bibitem{aassila} M.~Aassila, \emph{1000 challenges
 math\'ematiques, Alg\`ebre}, Ellipses, Paris, 2016, p.~110.
\bibitem{lean} F.~Gaudin, Formalization of OEIS A077754 in
 Lean~4 + mathlib (v4.35.0), archived at
 \url{https://doi.org/10.5281/zenodo.23879087}.
\end{thebibliography}

\end{document}
